\section{Method}
\label{sec:method}

\subsection{Problem Formulation}
\label{sec:formulation}

We formulate per-head KV cache quantization as a constrained optimization problem.
Let $h \in \{1,\ldots,H\}$ index the KV heads across all layers, $b_h \in \mathcal{B}$ the bit-width assigned to head $h$, and $\mathcal{B} = \{2, 4, 8, 16\}$ the set of hardware-aligned bit-widths.
Each head $h$ has a profiled distortion function $d_h(b)$, the PPL increase when head $h$ is quantized to $b$ bits.
Given a memory budget ratio $R$ (e.g., $R = 0.25$ for $4\times$ compression), we seek the allocation that minimizes total distortion:
\begin{equation}
\label{eq:objective}
\min_{\{b_h\}} \sum_{h=1}^{H} d_h(b_h) \quad \text{s.t.} \quad \frac{1}{H \cdot 16}\sum_{h=1}^{H} b_h \leq R.
\end{equation}

This is a multiple-choice knapsack problem (MCKP): each head selects one bit-width from $\mathcal{B}$ to minimize total distortion under a budget.
RateQuant~\cite{zuo2026ratequant} solves this for consecutive bit-widths $\{2,3,4,5,6\}$, where the polymatroid structure of the feasible set guarantees that greedy allocation is optimal (their Proposition 6).
Hardware-aligned bit-widths $\{2,4,8,16\}$ are non-consecutive; the gaps break this structure, and greedy is no longer guaranteed to match the MCKP optimum.
Moreover, as we show in \S\ref{sec:compounding}, the objective itself is not well-behaved: distortions from 2-bit heads compound non-additively, making unconstrained optimization unreliable.

\subsection{PPL-based Per-Head Profiling}
\label{sec:profiling}

We profile each KV head independently by quantizing it to each $b \in \mathcal{B}$ while keeping all other heads at bf16, and measuring the resulting PPL increase $d_h(b)$ on a calibration set of 64 sequences $\times$ 2048 tokens from WikiText-2 train.
This yields a distortion table of $H \times |\mathcal{B}|$ entries that serves as input to the solver.

The optimization in Eq.~\eqref{eq:objective} relies on an additive distortion assumption: the total PPL degradation from a mixed-precision allocation is approximately $\sum_h d_h(b_h)$.
This holds when per-head distortions are small and independent---a reasonable approximation for 4-bit and 8-bit quantization, where individual head distortions are typically below 0.1 PPL.
At 2-bit, however, individual distortions are orders of magnitude larger, and the additive assumption breaks down.

\subsection{Compounding Effect}
\label{sec:compounding}

When many heads are simultaneously quantized to 2-bit, per-head distortions compound non-additively.
The mechanism is straightforward: each head's quantized output feeds into the residual stream, and truncation errors from multiple 2-bit heads accumulate through subsequent layers.
When the number of 2-bit heads is small, these errors are absorbed by the remaining high-precision heads.
As the count grows, cumulative errors exceed the capacity of the residual stream to compensate, and PPL degrades super-linearly.

\begin{figure}[t]
\centering
\includegraphics[width=\columnwidth]{figures/fig2_compounding_curve.pdf}
\caption{PPL on Qwen2.5-7B ($4\times$ compression) as a function of $k$, the maximum number of 2-bit heads. Too few 2-bit heads waste budget on 4-bit; too many trigger compounding. The minimum occurs at $k{=}8$.}
\label{fig:compounding}
\end{figure}

Unconstrained solvers are blind to this effect because the additive objective underestimates the cost of large-scale 2-bit assignment.
On Qwen2.5-7B, unconstrained greedy assigns 68 heads to 2-bit and achieves 32.45 PPL; unconstrained MCKP assigns 75 heads and achieves 36.92---both far worse than uniform 4-bit (15.45).
The additive model predicts $\sim$11 PPL for these allocations, underestimating the actual degradation by $3\text{--}4\times$.

This compounding is specific to 2-bit quantization.
At 4-bit and above, per-head distortions are small enough that additivity holds even when all heads share the same bit-width (uniform 4-bit achieves 15.45 PPL, close to the sum-of-distortions prediction).
The qualitative difference arises because 2-bit quantization has far larger per-head distortion, saturating the error tolerance of the residual stream.

Figure~\ref{fig:compounding} shows the relationship between $k$ (the maximum number of 2-bit heads) and PPL on Qwen2.5-7B at $4\times$ compression.
The curve is U-shaped: at $k{=}0$ (uniform 4-bit), PPL is 15.45; allowing a few 2-bit heads reduces PPL by freeing budget for 8-bit heads on sensitive positions, reaching a minimum of 11.23 at $k{=}8$; beyond this point, compounding drives PPL upward, reaching 35.35 at $k{=}75$.

\subsection{Cardinality Constraint}
\label{sec:cardinality}

The compounding effect motivates a simple addition to Eq.~\eqref{eq:objective}: a cardinality constraint on 2-bit assignments.
\begin{equation}
\label{eq:constraint}
\sum_{h=1}^{H} \mathbb{1}[b_h = 2] \leq k.
\end{equation}

The constraint parameter $k$ is selected by a small-scale sweep on calibration data: solve the constrained problem for each candidate $k$, evaluate calibration PPL, and pick the value at the minimum of the U-shaped curve (Figure~\ref{fig:compounding}).
For Qwen2.5-7B with 112 KV heads, $k{=}8$ ($\sim$7\% of heads).

The constraint is a single linear inequality, compatible with any solver---greedy, MCKP, or integer linear programming---and adds negligible overhead to the optimization.

A natural question is whether Layer-0 heads require explicit protection, since attention heads in the first layer are known to carry disproportionate importance.
On Qwen2.5-7B, the profiler assigns very high distortion to Layer-0 heads (e.g., L0H0: $+$31.9\%, L0H2: $+$23.2\%, L0H3: $+$11.7\% PPL increase at 2-bit).
These values are large enough that no solver would assign 2-bit to Layer-0 under any reasonable budget.
We verify this empirically: running the solver with and without manual Layer-0 exclusion produces identical allocations across all 112 heads.
No manual exclusion rule is needed; the profiler signal is sufficient.

\subsection{Solver Equivalence}
\label{sec:solver_equiv}

\begin{table}[t]
\centering
\small
\caption{Solver comparison on Qwen2.5-7B at $4\times$ compression. Unconstrained solvers over-allocate 2-bit heads and fail catastrophically. The cardinality constraint ($k{\leq}8$) makes greedy and MCKP produce identical allocations.\footnotemark}
\label{tab:solver_equiv}
\begin{tabular}{llcccc}
\toprule
Solver & Constraint & PPL & \multicolumn{3}{c}{Bit distribution} \\
 & & & 2-bit & 4-bit & 8/16-bit \\
\midrule
Greedy & None & 32.45 & 68 & 30 & 14 \\
MCKP & None & 36.92 & 75 & 33 & 4 \\
Greedy & $k{\leq}8$ & 11.23 & 8 & 100 & 4 \\
MCKP & $k{\leq}8$ & 11.23 & 8 & 100 & 4 \\
Uniform & --- & 15.45 & 0 & 112 & 0 \\
\bottomrule
\end{tabular}
\end{table}
\footnotetext{Standard-resolution profiling; hires profiling achieves 10.91 (Table~\ref{tab:main_results}).}

With the cardinality constraint in place, greedy and MCKP produce identical results---not approximately, but exactly.
Table~\ref{tab:solver_equiv} shows that constrained greedy and constrained MCKP achieve the same PPL (11.23) with the same bit distribution on Qwen2.5-7B.
We verify this at the per-head level across three models:

\begin{itemize}[nosep]
\item \textbf{Qwen2.5-7B:} 112/112 heads receive identical bit-widths.
\item \textbf{Llama-3.1-8B:} 256/256 heads identical.
\item \textbf{Mistral-7B-v0.3:} 255/256 heads identical. The single difference is a zero-distortion tie between two heads at 4-bit and 8-bit, with no PPL impact.
\end{itemize}

The constraint narrows the feasible set enough that greedy's locally optimal decisions are globally optimal.
Without it, greedy cascades into excessive 2-bit allocation: each additional 2-bit head looks locally beneficial (freeing budget) but contributes to the compounding effect that degrades global PPL.
The cardinality constraint breaks this cascade by capping the number of 2-bit heads, aligning local and global optima.

The practical implication is that practitioners can use the simpler $O(H \log H)$ greedy algorithm instead of MCKP without sacrificing optimality.
